% === §2.2 Pertinencia de la propuesta ===
% Redactado por el Redactor — Fase 2

% --- Párrafo 1: Cifras socioeconómicas/tecnológicas, ODS y políticas; articulación con criterio e ---

La pertinencia del proyecto se sustenta en indicadores convergentes del contexto educativo, tecnológico y productivo colombiano. El Ministerio de Educación Nacional reporta tasas de deserción interanual cercanas al 45\,\% en programas de ingeniería y ciencias exactas, mientras el PND 2022--2026 y el Documento CONPES 4075 de 2022 advierten un déficit estructural de talento humano en tecnologías habilitadoras de la Industria 4.0, estimado en más de 60\,000 profesionales de TI no cubiertos. A escala global, la OCDE (2023) estima que un 40\,\% de las competencias técnicas requeridas por el sector productivo se habrán transformado hacia 2027, y la UNESCO, en su Estrategia de Educación 2030, prioriza la integración de la IA como palanca de calidad y equidad educativa \cite{unesco2023aied}. El proyecto se alinea directamente con el ODS~4 (Educación de calidad e inclusiva), el ODS~9 (Industria, innovación e infraestructura) y el ODS~10 (Reducción de desigualdades), al democratizar el acceso a un tutor inteligente de frontera para estudiantes de instituciones públicas; con el Plan Global de Desarrollo UNAL 2025--2027, en su línea de apropiación social del conocimiento; y con la Misión de Sabios 2019,\,eje de la cuarta transformación. En el marco del criterio~\textbf{e} de la Convocatoria Nacional Alianzas SIUN 2025--2027 (aportes a necesidades territoriales, nacionales y globales, 20 pts), el ecosistema agéntico aporta a nivel \emph{territorial} mediante un piloto en la Sede Manizales y su área de influencia cafetera; a nivel \emph{nacional} al integrarse al Sistema de Investigación de la UNAL (SIUN) en carácter de alianza; y a nivel \emph{global} al generar conocimiento transferible sobre IA agéntica educativa publicable en revistas Q1/Q2.

% --- Párrafo 2: Justificación técnica del uso de IA (predictiva, generativa, agéntica) frente a métodos tradicionales; criterio a (calidad) y b (articulación SIUN) ---

 Desde el punto de vista técnico, el uso combinado de IA predictiva, generativa y agéntica ofrece ventajas medibles frente a los métodos tradicionales de tutoría y evaluación. La IA \emph{predictiva} permite modelar el estado cognitivo del aprendiz a partir de trazas multimodales (interacción con código, simulaciones y argumentación textual), superando los indicadores superficiales de los sistemas de gestión del aprendizaje convencionales \cite{hooshyar2024systematic}. La IA \emph{generativa}, soportada en grandes modelos de lenguaje, habilita la síntesis de problemas auténticos y la retroalimentación en lenguaje natural de alta calidad en dominios ingenieriles donde la elaboración manual de ítems es prohibitiva \cite{wang2024survey, park2023generative}. La IA \emph{agéntica}, finalmente, aporta orquestación deliberada: una arquitectura multiagente con roles pedagógicamente fundados coordina diagnóstico, tutorización, simulación y evaluación, manteniendo memoria de trayectoria y usando protocolos como el Model Context Protocol (MCP) para integrar herramientas externas de manera segura \cite{xi2023rise}. Mientras un sistema tutor inteligente basado en reglas requiere años de ingeniería del conocimiento por dominio, la arquitectura agéntica propuesta reconfigura sus estrategias a partir del modelo del aprendiz, escalando a múltiples asignaturas con un costo marginal bajo. Esta integración constituye una innovación metodológica comprobable (criterio~\textbf{a} de la convocatoria, 30 pts) y articula de forma directa las líneas de investigación del GCPDS (categoría A1 Minciencias) en procesamiento de señales, aprendizaje automático y educación en ingeniería con el LabIA y el SIUN (criterio~\textbf{b}, 25 pts), consolidando una alianza intra-SIUN con capacidades complementarias: modelado del aprendiz, arquitecturas multiagente y validación empírica en aula.

% --- Párrafo 3: Producto tangible esperado, TRL 6/7 y transferencia; criterio c (articulación externa) y d (formación de estudiantes) ---

El producto tangible esperado es un \emph{ecosistema de software de agentes autónomos para educación en ingeniería}, desplegado tanto en infraestructura local (estaciones Apple Silicon del LabIA) como en la nube (VPS), con interfaz conversacional, integración mediante MCP a herramientas de simulación eléctrica/electrónica y de programación, y un panel de analítica del aprendizaje para docentes. Al cierre del proyecto el ecosistema alcanzará un \textbf{TRL~6} demostrado en un entorno relevante (asignaturas reales de pregrado y posgrado de la UNAL) y se proyectará hacia el \textbf{TRL~7} en un entorno operativo fomentando la adopción por parte de grupos de investigación aliados y del sector productivo —empresas de software y formación técnica del eje cafetero—, lo que da cumplimiento al criterio~\textbf{c} (articulación con actores externos, 5 pts). La transferencia tecnológica se materializa mediante código abierto bajo licencia permisiva, un registro de software ante la Dirección Nacional de Derechos de Autor de Colombia, y la documentación de despliegue que permita a terceros instalar y operar el ecosistema. La pertinencia formativa es central en la propuesta (criterio~\textbf{d}, 20 pts): el proyecto servirá como vehículo de formación en investigación para al menos seis estudiantes de pregrado (Ingeniería Electrónica, Ingeniería de Sistemas, Ciencias de la Computación) y tres de posgrado (Maestría y Doctorado en Ingeniería Automática), en roles de desarrollo agéntico, curaduría de datos, evaluación empírica y redacción científica, con mentoría directa de los investigadores del GCPDS—LabIA—UNAL. En síntesis, el ecosistema conecta la frontera de la IA agéntica con sus beneficiarios directos —estudiantes, docentes y empresas— rindiendo evidencia cuantitativa de impacto y un artefacto tecnológico transferible, condición indispensable para que la apropiación del conocimiento supere el umbral del laboratorio.